\section*{Proof of Theorem 1 (original Section 2)}
The three players in our protocol are called $P_x$, $P_y$ and $P_z$. The inputs that they get to see are $(y,z)$, $(x,z)$ and $(x,y)$ respectively.
Here is a similar problem in the realm of vector addition. Given integers $q,d>1$, define $g=g_{q,d}(\alpha,\beta,\gamma)$ to be 1 if $\alpha+\beta=\gamma$ and 0 otherwise. Here $\alpha,\beta\in[q]^d$, $\gamma\in[2q]^d$ and addition is vector addition in $\mathbb R^d$. The following one-round protocol [3] for $g$ is correct because the inequality $\norm{2\alpha-\gamma}^2+\norm{2\beta-\gamma}^2\ge2\norm{\alpha-\beta}^2$ holds always and is an equality iff $\gamma=\alpha+\beta$.
\paragraph{Protocol 1. A protocol for $g_{q,d}$.}
\begin{enumerate}
\item $P_z$ computes $\norm{\alpha-\beta}_2^2$, and writes the result on the board.
\item $P_y$ writes 1 iff $\norm{\alpha-\beta}_2^2=\norm{2\alpha-\gamma}_2^2$.
\item $P_x$ writes 1 iff $\norm{\alpha-\beta}_2^2=\norm{2\beta-\gamma}_2^2$.
\end{enumerate}
The cost of this protocol is $2+\log dq^2$.
The above is an efficient method to decide high-dimensional vector addition, but our objective is to decide the integer addition relation $X+Y+Z=N$. We let $x=X$, $y=Y$ and $z=N-Z$, so the relation we need to consider is $x+y=z$.
Our protocol to decide whether $x+y=z$ builds on the protocol for $g_{q,d}$. It is the issue of carry bits in integer addition that makes this decision problem harder. The integers $q,d>1$ are chosen so that
\begin{equation}2qN>q^d\ge2N.\tag{3}\label{eq:range}\end{equation}
the specific choice is made below so as to minimize the cost of the protocol.
We denote by $w_q$ the vector that corresponds to the base $q$ representation of the integer $w$. As usual, $e_i$ is the $d$-dimensional vector with 1 in the $i$-th coordinate and zeros elsewhere.
Let $C(x,y)\in\{0,1\}^d$ be the carry vector when $x$ and $y$ are added in base $q$. The relation $x+y=z$ among integers is equivalent to the vector relation
\[x_q+y_q=\zeta,\]
where the $i$-th coordinate of $\zeta$ is
\[\zeta_i=z_i+q\cdot C(x,y)_i-C(x,y)_{i-1}\]
(Here $C(x,y)_0=0$).
The protocol from [3] now suggests itself: $P_z$ posts $C(x,y)$, and Protocol 1 is used to decide the relation $x_q+y_q=\zeta$. This yields again the estimate (1) of the original paper, $2\sqrt2\sqrt{\log N}+o(\sqrt{\log N})$.
The alternative approach that we adopt here considers instead the equivalent vector relation
\[x_q+\eta=z_q\]
where
\[\eta=(x+y)_q-x_q.\]
Concretely, the $i$-th coordinate of $\eta$ is:
\[\eta_i=y_i-q\cdot C(x,y)_i+C(x,y)_{i-1}.\]
In order to run Protocol 1, $P_z$ needs to know $\eta$ and $x_q$, which he does. The situation with $P_y$ is even simpler, since he needs to know $x_q$ and $z_q$ which are his inputs. The only difficulty is with $P_x$ who needs to know $z_q$ (which he does) and $\eta$. The latter is not part of his input and $P_z$ fills in the missing information for him.
The obvious solution is for $P_z$ to reveal $C$ to $P_x$ using $d$ bits of information. However, we can save communication by exploiting the fact that $P_x$ and $P_z$ share some information, i.e., they both know $y$ for every $y\ne0$.
By a standard argument in this area which we detail below (Proposition 2), a protocol that works for typical pairs $x,y$ can be easily modified to work in all cases. So, let us pick $x$ and $y$ uniformly at random from among the $d$-digit numbers in base $q$ and think of $C$, the vector of carry bits as a random variable on this probability space. The number of bits that $P_z$ needs to post so that $P_x$ gets to know $C$, and therefore know $\eta$, is $H(C\mid y)$, the entropy of $C$ given $y$. The gain is clear, since $H(C)\ge H(C\mid y)$.
It remains to estimate $H(C\mid y)$. Fix some integers $s\ge t\ge0$, and let $X$ be the random variable that is a uniformly sampled subset of $[s]$ of cardinality $\ge t$. It is easily verified that $H(X)=(1+o_s(1))\cdot s\cdot h(t/s)$, where $h(\cdot)$ is the univariate entropy function. The entropy of $X$ is the same also if we sample subsets of $[s]$ of cardinality $<t$. Let $r$ be an integer in the range $d\gg r\gg1$, e.g., $r\approx\sqrt d$. For $j=1,\ldots,r$, let
\[S_j=\{i\mid \tfrac{qj}{r}>y_i\ge\tfrac{q(j-1)}{r}\},\]
where $q>y_i\ge0$ is the $i$-th digit of $y$. A carry occurs in digit $i\in S_j$ only if $x_i>\frac{q(r-j)}r$, where $x_i$ is the $i$-th digit of $x$. Then
\[H(C\mid y)\le(1+o_r(1))\sum_{j=1}^r\frac{|S_j|}{d}h(\tfrac jr).\]
Since $y$ is chosen at random, $|S_j|\le(1+o_r(1))\frac dr$, and so
\[H(C\mid y)\le(1+o_r(1))\sum_{j=1}^r\frac1r h(\tfrac jr).\]
The limit of this expression as $r\to\infty$ is
\[\lambda=\int_0^1h(u)\,du=\frac{\log e}{2}=0.721\ldots\]
It is left to optimize on $q$ and $d$. The complexity of our protocol is
\[\lambda d+\log dq^2+2,\]
where recall that $2qN>q^d\ge2N$. It is not hard to verify that choosing
\begin{equation}d=\sqrt{\frac2\lambda\log2N}\qquad q=2^{\sqrt{\frac\lambda2\log2N}},\tag{4}\label{eq:choice}\end{equation}
we get a protocol with complexity
\[2\sqrt{2\lambda\log N}+o(\sqrt{\log N}),\]
and this is asymptotically optimal in our setting.
To sum up, here is the protocol which proves Theorem 1:
\paragraph{Protocol 2. A protocol for exactly-$N$, for typical pairs $x,y$.}
For $d,q$ as in Equation (4)
\begin{enumerate}
\item $P_z$ publishes the vector $\eta=(x+y)_q-x_q$ in a way that $P_x$ can read.
\item The players run protocol 1 for $g_{q,d}$ on inputs $x_q,\eta,z_q$. That is:
\begin{enumerate}
\item $P_z$ writes $\norm{\eta-x_q}_2^2$ on the board
\item $P_y$ writes 1 iff $\norm{\eta-x_q}_2^2=\norm{2x_q-z_q}_2^2$.
\item $P_x$ writes 1 iff $\norm{\eta-x_q}_2^2=\norm{2\eta-z_q}_2^2$.
\end{enumerate}
\end{enumerate}
\begin{proposition}
Let $P$ be an NOF protocol for the exactly-$N$ that works correctly for an $\Omega(1)$-fraction of the input pairs $x,y$ (and every $z$) with communication complexity $\Phi(N)$. Then there is an NOF protocol that works for all inputs with communication complexity $\Phi(N)+O(\log\log N)$.
\end{proposition}
\begin{proof}
Let $S\subseteq[N]\times[N]$ be the set of input pairs $x,y$ on which $P$ succeeds. We claim that there is a collection $F$ of $O(\log N)$ vectors $\Delta\in[N]\times[N]$ such that
\[\bigcup_{\Delta\in F}(S+\Delta)\supseteq[N]\times[N].\]
In the modified protocol, $P_z$ sees $x,y$ and announces the index of some $\Delta=(\Delta_1,\Delta_2)\in F$ for which $(x-\Delta_1,y-\Delta_2)\in S$. Then the players run Protocol 2 with inputs $(x-\Delta_1,y-\Delta_2,z-\Delta_1-\Delta_2)$.
The construction of $F$ uses a standard fact about the set-cover problem. For a family of finite sets $\mathcal X\subseteq2^\Omega$ we denote by $c(\mathcal X)$ the least number of members in $\mathcal X$ whose union is $\Omega$. Also $c^*(\mathcal X)$ is the minimum cost of a fractional cover. Namely,
\[c^*(\mathcal X)=\min\sum_{\mathcal X}\omega_X,\quad\text{where }\omega_X\ge0\text{ for every }X\in\mathcal X\text{ and }\sum_{x\in X}\omega_X\ge1\text{ for every }x\in\Omega.\]
Then
\[c(\mathcal X)\le\log(|\Omega|)\cdot c^*(\mathcal X)\]
(e.g., Lov\'asz [15]) and actually the greedy algorithm yields a set cover that meets this bound.
In our case $\Omega=[N]\times[N]$, and
\[\mathcal X=\{(S+\Delta)\cap([N]\times[N])\mid\Delta\in[-N,N]\times[-N,N]\}.\]
It is easily verified that the weights $\omega_x=\frac{10}{N^2}$ constitute a fractional cover, so that $c^*(\mathcal X)\le40$ and hence $c(\mathcal X)\le80\log N$, as claimed.
\end{proof}
